% !TeX root = ../../main.tex
\subsection{Boundary exponentials and forgetting markings}
\label{subsec:s02:04}

We first expand the boundary exponential using the normal line of a
node.  We then recall the forgetting formulas and invert their cycle
sum.  A degree lemma isolates the rational-tail corrections, so that
the final exponential specialization can also be used with existing
cotangent and kappa insertions.

For formal variables $x,y$ and a scalar parameter $a$, define the
symmetric edge kernel
\begin{equation}
\label{eq:s02:edge-kernel}
K_a(x,y)=\sum_{r=1}^{\infty}\frac{a^r(-x-y)^{r-1}}{r!}
=\frac{1-e^{-a(x+y)}}{x+y}
=\int_0^a e^{-t(x+y)}\,\dd t.
\end{equation}
Here $r$ runs over the positive integers, and the integral is formal
termwise integration in $t$.  The numerator of the fraction is
divisible by $x+y$, so no inversion of a cotangent class is involved.

\begin{prp}[Boundary exponential]
\label{prp:s02:boundary-exponential}
Let $a,c$ be scalar parameters, and let $d_1,\ldots,d_n$ be
nonnegative integers.  Then
\begin{equation}
\label{eq:s02:boundary-exponential}
\begin{split}
e^{a\delta+c\kappa_1}\prod_{i=1}^n\psi_i^{d_i}
={}&\sum_{\Gamma\in\mathcal G_{g,n}}
\frac{1}{|\operatorname{Aut}\Gamma|}\,
(\xi_\Gamma)_*\left[
\prod_{v\in V(\Gamma)}e^{c\kappa_1^{(v)}}
\prod_{i=1}^n\psi_i^{d_i}
\prod_{\{f,f'\}\in E(\Gamma)}K_a(\psi_f,\psi_{f'})
\right].
\end{split}
\end{equation}
All classes in brackets lie on $\MM_\Gamma$; a leg class is on its
incident factor.  Both sides are truncated by codimension.  The graph
with no edges contributes the original vertex expression.
\end{prp}

This is the constant-edge specialization of Andersen et al.,
\cite[Lemma 3.10, (3.18)--(3.22)]{ABCDGLW23}, using the numbering of
arXiv:1905.10352v3.  Their edge series is the constant $a$, and their
vertex series is linear with coefficient $c$.  Their sum of one-edge
pushforwards weighted by $1/|\operatorname{Aut}\Gamma|$ is our
boundary class $\delta$.  The cited statement is in cohomology; the
same normal-bundle intersection proof gives the present Chow identity,
as follows.

\begin{prf}
The smoothing parameter of the node for $\{f,f'\}$ belongs to the
tensor product of its two tangent lines.  Its normal line therefore
has first Chern class $-\psi_f-\psi_{f'}$.  In a product of boundary
divisors, repeating this node $r$ times contributes
$(-\psi_f-\psi_{f'})^{r-1}$ after the gluing pushforward.
Distinct nodes intersect as the corresponding boundary branches.
Grouping the terms of the exponential by their graph of distinct
nodes gives at each edge the series
$\sum_{r\geq1}a^r(-\psi_f-\psi_{f'})^{r-1}/r!$, which is
\eqref{eq:s02:edge-kernel}.  The factorial of the exponential cancels
the orderings of the chosen divisor factors; passing from distinguished
graph labels to graph isomorphism classes leaves
$1/|\operatorname{Aut}\Gamma|$.
These are algebraic self-intersection and pushforward calculations,
valid in rational Chow on the moduli stack.  Finally, the projection
formula and \eqref{eq:s02:tautological-gluing} supply the vertex
exponentials and leg factors in \eqref{eq:s02:boundary-exponential}.
\end{prf}

For example, the one-loop graph with vertex type $(0,3)$ and one leg
has two automorphisms, obtained by exchanging its loop flags.  In the
bare boundary exponential on $\MM_{1,1}$ its term is
$(a/2)(\xi_\Gamma)_*1$.  A separating edge on $\MM_{0,4}$ with leg
sets $\{1,2\}$ and $\{3,4\}$ has no nontrivial automorphism and
contributes $a(\xi_\Gamma)_*1$.  More generally, interchanging the
two sides of a separating edge is an automorphism only when it
preserves their genera and every leg label; in the unmarked equal-genus
case it contributes a factor $1/2$.  Summing over ordered separating
splittings instead requires a uniform prefactor $1/2$, since unequal
sides then occur twice.  These conventions agree with
$\delta=\frac12\xi_*1$ in \eqref{eq:s02:ordered-boundary}.

For an integer $m\geq1$, write
$\pi_m=\pi_{n+m,n}:\MM_{g,n+m}\to\MM_{g,n}$, with
$\pi_0$ the identity.  In a pullback formula, classes inside
$\pi_m^*$ are on the base, and classes outside it are intrinsic to
the source.  For $1\leq i\leq n$, let
$D_{0,\{i,n+1\}}$ be the divisor whose rational tail carries exactly
the markings $i,n+1$.  The one-point identities used below are
\begin{equation}
\label{eq:s02:one-point-comparison}
\pi_1^*\psi_i=\psi_i-D_{0,\{i,n+1\}},\qquad
\pi_1^*\kappa_a=\kappa_a-\psi_{n+1}^a
\quad(1\leq i\leq n,\ a\geq0).
\end{equation}
The first comparison follows from the
cotangent line under contraction of this three-pointed tail; its
vanishing along that divisor has multiplicity one.  The second is the
logarithmic-kappa comparison in \cite[(2.10)]{KMZ96}, whose classes
$\omega_n(a)$ are our $\kappa_a$.

For its algebraic justification, consider the universal curve over
$\MM_{g,n+1}$ and its stabilization to the pulled-back universal
curve over $\MM_{g,n}$.  On the former, let $D$ be the new section
divisor and let $\mathcal L$ be the pullback of the old logarithmic
dualizing line.  The contracted rational components have degree zero
for the dualizing line with only the old markings; stabilization
therefore identifies that line with $\mathcal L$.  The new
logarithmic dualizing line is $\mathcal L(D)$.  On $D$, the class
$c_1(\mathcal L)$ restricts to $\psi_{n+1}$, while the section's
self-intersection is $-\psi_{n+1}$.  The projection formula and flat
base change identify the pushforward of $c_1(\mathcal L)^{a+1}$
with $\pi_1^*\kappa_a$.  Expanding $c_1(\mathcal L(D))^{a+1}$ adds
\[
\psi_{n+1}^a\sum_{k=1}^{a+1}(-1)^{k-1}\binom{a+1}{k}
=\psi_{n+1}^a.
\]
This proves the kappa comparison in Chow without appealing to
injectivity of the cycle map.

\begin{lmm}[Rational-tail degree]
\label{lmm:s02:rational-tail-degree}
Fix $m\geq1$ and integers $r_1,\ldots,r_m\geq1$, and set
$W=\prod_{k=1}^m\psi_{n+k}^{r_k}$ on $\MM_{g,n+m}$.  Then
\begin{equation}
\label{eq:s02:tail-comparisons}
\begin{split}
W\pi_m^*\psi_i&=W\psi_i\quad(1\leq i\leq n),\\
W\pi_m^*\kappa_1
&=W\left(\kappa_1-\sum_{k=1}^m\psi_{n+k}\right).
\end{split}
\end{equation}
The same replacements hold in any polynomial or codimension-truncated
formal series in the old cotangent classes and $\kappa_1$, after
multiplication by $W$.
\end{lmm}

\begin{prf}
Before multiplication by $W$, iterating
\eqref{eq:s02:one-point-comparison} expresses the comparison errors
in \eqref{eq:s02:tail-comparisons} as linear combinations of
rational-tail divisors contracted by $\pi_m$.  For an old cotangent
class these tails carry that old marking; the remaining boundary
terms in the kappa comparison carry only new markings.  Pulling back a comparison
divisor at an intermediate forgetting step produces divisors of this
same kind.  Comparing an intermediate new cotangent class with its
class on the final source also produces only these divisors; this
gives the surviving sum of new cotangent classes in the second line.

Suppose one such tail carries $k$ new markings and $\epsilon$ old
markings, where $k\geq1$ and $\epsilon\in\{0,1\}$.  Its factor is
$\MM_{0,k+\epsilon+1}$.  Its dimension is $k-1$ with one old
marking and $k-2$ with none; only stable factors occur.  The
restriction of $W$ to it contains
new cotangent powers of total degree at least $k$, so is zero.
By the projection formula $W$ annihilates the class of the divisor
itself, and hence the ideal generated by all these comparison
divisors.  The first two equalities follow.  The polynomial statement
then follows by multiplication in the quotient by that ideal, and the
formal-series statement by codimension truncation.
\end{prf}

In particular, for scalar parameters $c,t_1,\ldots,t_n$ the lemma gives
\begin{equation}
\label{eq:s02:tail-exponential}
W\pi_m^*e^{c\kappa_1+\sum_{i=1}^nt_i\psi_i}
=W e^{c\kappa_1-c\sum_{k=1}^m\psi_{n+k}+\sum_{i=1}^nt_i\psi_i}.
\end{equation}
The term $-c\sum_k\psi_{n+k}$ remains.  This statement requires a
positive cotangent power at every forgotten marking, as in the lemma.

Let $[m]=\{1,\ldots,m\}$ and let $\mathfrak S_m$ be its permutation
group.  For $\sigma\in\mathfrak S_m$, denote by
$\operatorname{Cyc}(\sigma)$ its set of cycles, viewed as subsets of
$[m]$.  A set partition $\mathcal P$ of $[m]$ consists of unordered
nonempty blocks.  When writing its blocks as $C_1,\ldots,C_k$, with
$k=|\mathcal P|$, any ordering may be chosen.

\begin{lmm}[Forgetting and partition inversion]
\label{lmm:s02:forgetting-inversion}
For $m\geq1$ and nonnegative integers $a_1,\ldots,a_m$,
\begin{equation}
\label{eq:s02:forgetting-permutations}
(\pi_m)_*\prod_{i=1}^m\psi_{n+i}^{a_i+1}
=\sum_{\sigma\in\mathfrak S_m}
\prod_{C\in\operatorname{Cyc}(\sigma)}\kappa_{\sum_{i\in C}a_i}.
\end{equation}
Its inverse is
\begin{equation}
\label{eq:s02:partition-inversion}
\prod_{i=1}^m\kappa_{a_i}
=\sum_{\mathcal P\text{ a partition of }[m]}
(-1)^{m-|\mathcal P|}(\pi_{|\mathcal P|})_*
\prod_{k=1}^{|\mathcal P|}
\psi_{n+k}^{\,1+\sum_{i\in C_k}a_i}.
\end{equation}
The pushforward is independent of the chosen ordering of the blocks.
\end{lmm}

These are Kaufmann--Manin--Zagier's formulas
\cite[Theorem 2.2, (2.6), and Lemma 2.2.1, (2.8)]{KMZ96}, with
numbering from arXiv:alg-geom/9604001v1.  They credit the geometric
cycle formula to Arbarello--Cornalba and note its earlier unpublished
derivation by Faber--Zagier.  Their
$\omega_n(a_1,\ldots,a_m)$ in (2.4) is the left side of
\eqref{eq:s02:forgetting-permutations}, and $\omega_n(a)=\kappa_a$.
Their inverse sums over ordered blocks with a factor $1/k!$; summing
over unordered partitions gives \eqref{eq:s02:partition-inversion}.

\begin{prf}
The geometric inputs are one-point forgetting and the logarithmic
comparison \eqref{eq:s02:one-point-comparison}.  The old cotangent
pullback corrections are killed by the positive power at the point
being forgotten.  Thus their push-pull computation, recorded in
\cite[(2.9)--(2.11)]{KMZ96}, applies in Chow as well as cohomology.
It gives a recursion for the left side of
\eqref{eq:s02:forgetting-permutations}: the last label either supplies
its own kappa factor, or its degree is added at one of the other
labels.  These alternatives correspond exactly to the last label
being a singleton cycle or being inserted after a specified label in
a cycle.  Starting with $m=1$ proves the permutation formula.

To verify the inverse, substitute the permutation formula into its
right side and group by the resulting blocks of original labels.  A
block formed from $k$ intermediate blocks acquires the cycle factor
$(k-1)!$.  The signed generating series for one intermediate nonempty
block of labelled elements is $1-e^{-z}$, for a formal variable $z$.
Consequently, the generating series for each resulting block is
\[
\sum_{k=1}^{\infty}\frac{(1-e^{-z})^k}{k}
=-\log(e^{-z})=z.
\]
Here division by $k!$ for unordered intermediate blocks combines with
$(k-1)!$ to give $1/k$.  Only singleton resulting blocks survive,
with coefficient one, proving \eqref{eq:s02:partition-inversion}.
\end{prf}

For instance, for nonnegative integers $a,c$,
\begin{equation}
\label{eq:s02:two-kappas}
\kappa_a\kappa_c
=(\pi_2)_*(\psi_{n+1}^{a+1}\psi_{n+2}^{c+1})
-(\pi_1)_*\psi_{n+1}^{a+c+1}.
\end{equation}
The inverse has no block factorials; the factorials $(|C|-1)!$ occur
only when cycles with the same underlying block are counted in the
forward permutation sum.

\begin{crl}[Exponential kappa specialization]
\label{crl:s02:kappa-exponential}
For a scalar parameter $s$,
\begin{equation}
\label{eq:s02:kappa-exponential}
e^{s\kappa_1}
=\sum_{m=0}^{\infty}\frac{1}{m!}(\pi_m)_*
\prod_{i=1}^m\psi_{n+i}(1-e^{-s\psi_{n+i}}).
\end{equation}
The empty product is one.  At $s=1$ this is the desired expansion of
$e^{\kappa_1}$ by additional cotangent insertions.  At fixed $(g,n)$
the right side is finite after codimension truncation.
\end{crl}

\begin{prf}
Expand $e^{s\kappa_1}$ and apply
\eqref{eq:s02:partition-inversion} to each power of $\kappa_1$.
If a partition has $m$ blocks of sizes $r_1,\ldots,r_m\geq1$, its
contribution, after cancelling the factorial for the original labels,
has coefficient
\[
\frac{1}{m!}\prod_{i=1}^m
\frac{(-1)^{r_i-1}s^{r_i}}{r_i!}.
\]
The corresponding new marking carries $\psi_{n+i}^{r_i+1}$.
For $1\leq i\leq m$, summing its block size gives
\[
\sum_{r=1}^{\infty}\frac{(-1)^{r-1}s^r}{r!}\psi_{n+i}^{r+1}
=\psi_{n+i}(1-e^{-s\psi_{n+i}}).
\]
This proves \eqref{eq:s02:kappa-exponential}, including all block
and marking factorials.  Each new factor has minimum degree two;
forgetting $m$ markings lowers codimension by $m$, so its pushforward
has minimum degree $m$.  Thus terms with $m>3g-3+n$ vanish.
\end{prf}

The specialization can be multiplied by a Hodge insertion and old
cotangent factors using the projection formula.  Hodge classes pull
back unchanged, and Lemma~\ref{lmm:s02:rational-tail-degree} replaces
the pulled-back old cotangent classes by those on the source, because
each new factor has degree at least two.  For an additional kappa
background the replacement is instead exactly
\eqref{eq:s02:tail-exponential}.  These are finite-genus identities;
they do not provide a uniform bound for subsequent genus-dependent
sums.
